%% Generated by tools/md2tex.js from paper/software.md. Do not edit by hand: edit
%% software.md, then run  node tools/md2tex.js software  and recompile.
\documentclass[11pt]{article}
\usepackage[T1]{fontenc}
\usepackage{lmodern}
\usepackage[utf8]{inputenc}
\usepackage[margin=1in]{geometry}
\usepackage{booktabs,graphicx,amsmath,amssymb}
\graphicspath{{../}}
\usepackage{tabularx}
\setlength{\emergencystretch}{3em}
\usepackage[numbers,sort&compress]{natbib}
\usepackage{threeparttablex}
\usepackage[british]{babel}
\usepackage[hidelinks]{hyperref}
\setlength{\bibsep}{2pt}
% [S#] is a provenance tag. Each one resolves to a row of the appendix table in
% software.md, which paper/tex_audit-software.md re-checks file by file. It is rendered as a
% small superscript rather than dropped, so an unresolved tag stays visible.
\newcommand{\prov}[1]{\textsuperscript{\,#1}}
\widowpenalty=10000
\clubpenalty=10000
\renewcommand{\topfraction}{0.9}\renewcommand{\bottomfraction}{0.8}
\renewcommand{\textfraction}{0.07}\renewcommand{\floatpagefraction}{0.75}
\setcounter{topnumber}{3}\setcounter{totalnumber}{4}
\tolerance=2000
\emergencystretch=2em
\renewcommand{\tabularxcolumn}[1]{>{\raggedright\arraybackslash}p{#1}}
\title{nlcsplit: fragment and pair resolution of the VV10 nonlocal-correlation energy}
\author{Angran Xia\thanks{Corresponding author; e-mail in the byline.} \\ \small \textit{School of Chemistry \& Chemical Engineering, Shaanxi Normal University, Xi'an 710119, China} \\ \small \texttt{haangyeon@snnu.edu.cn} \and Changwei Wang \\ \small \textit{School of Chemistry \& Chemical Engineering, Shaanxi Normal University, Xi'an 710119, China}}
\date{Draft --- 29 September 2026\\ \small \textit{ORCIDs PENDING, author-supplied}}
\begin{document}
\maketitle
% byline/affiliation already typeset by uthor

\section{Abstract}

% byline/affiliation already typeset by uthor

\section{Motivation and statement of need}

Functionals that pair a range-separated hybrid with the VV10 nonlocal correlation term - $\omega$B97X-V \citep{mardirossian2014wb97xv} and $\omega$B97M-V \citep{mardirossian2016wb97mv} in particular - are a routine choice for non-covalent interactions. Their nonlocal contribution,

    E\_NL = $\frac{1}{2}$ $\int$$\int$ $\rho$(r) $\rho$(r') $\Phi$(r, r') d$^{3}$r d$^{3}$r',

is a two-centre sum, but what users get back is one number per grid point, or one total. When an interaction energy needs to be interpreted per fragment, the available options stop at the semi-local part: electrostatics, exchange, polarisation and charge transfer can be resolved by EDA schemes, whereas the nonlocal term is carried as an undivided block. The closest published treatment computes exactly that block as a whole and asks whether it is accurate \citep{dasgupta2023}; the closest published explicit evaluation, JuNoLo \citep{lazic2010junolo}, sums the double sum directly but targets an improved total energy and explicitly recommends itself for bonded systems "where one cannot tell apart fragments". Neither resolves the term into fragment pairs, which is what the literature screen of \prov{S1} was run to find - and that screen's own record says the open-index \emph{subject} search behind it is unreliable, so this is a screen, not an enumeration. That gap is what this package fills, and it is a capability gap rather than an accuracy gap: on the same grid and density our sum agrees with the native PySCF kernel to \textasciitilde{}1e-8 Ha \prov{S9}.

\section{Method overview}

\textbf{Kernel and convention.} \texttt{w0 = sqrt(C|$\nabla$$\rho$|$^{4}$/\allowbreak{}$\rho$$^{4}$ + 4$\pi$$\rho$/\allowbreak{}3)}, \texttt{$\kappa$ = b(3$\pi$/2)($\rho$/9$\pi$)\textasciicircum{}(1/6)}, \texttt{g = w0(r\_i)R$^{2}$ + $\kappa$(r\_i)} with same-centre pairing, \texttt{$\Phi$ = -3/(2 g g' (g + g'))}, and \texttt{E\_NLC = $\frac{1}{2}$$\iint$$\rho$$\rho$$\Phi$ + $\beta$$\int$$\rho$} with \texttt{$\beta$ = (1/32)(3/b$^{2}$)\textasciicircum{}(3/4)}. Parameters are read from PySCF \texttt{nlc\_coeff} rather than hard-coded. Same-centre pairing is the convention that reproduces the native kernel; the cross-pairing variant \texttt{g = w0(r\_j)R$^{2}$ + $\kappa$(r\_i)} is retained as a definition-sensitivity control and raises monomer energies by +0.1379 (water), +0.4182 (argon) and +0.1346 (methane) kcal/mol at grid level 1, and +0.9057 for benzene at level 0 \prov{S5}. The offset is not a constant and it moves with the grid - the level-0 benzene figure is not the level-1 one - so these four numbers are one batch at the levels the producing script itself selects, and must not be mixed with another batch. \prov{S5} now points at an archived run carrying all four, which it did not: an earlier version of this sentence gave benzene as 0.914 kcal/mol and labelled the whole list "at grid level 1", and neither value nor label was checkable, because the only carrier had been a gitignored working log. The rerun that closed the gap reproduced water and argon to the printed digit and gave benzene +0.9057 at level 0 - which is what the companion text already said, so the two papers had been disagreeing with each other about one number while citing the same source row - and added methane, which neither paper had mentioned.

\textbf{Blocks and additivity.} \texttt{pair\_energy(A, B)} returns the half block $\frac{1}{2}$ $\Sigma$\_{i$\in$A} $\Sigma$\_{j$\in$B}; the inter-fragment term is reported as twice that, so \texttt{E\_\allowbreak{}nl = $\Sigma$\_\allowbreak{}A intra[A] + $\Sigma$\_\allowbreak{}{A<B} inter[A,B]} holds as an identity. Measured residual: 1.73e-18 to 6.94e-18 Ha, over the five partitions and two repeated schemes of the archived level-1 water-dimer run \prov{S2}.

\textbf{Why decomposing only the double sum is exact.} $\beta$ is a constant and $\int$$\rho$ is the electron number, so the local term $\beta$$\int$$\rho$ = $\beta$N cancels identically from any superposition difference. Inter-fragment decomposition therefore needs no convention for splitting a local energy density: it is a property of the double sum alone \prov{S3}. The flip side is reported as a limitation rather than hidden - the absolute local term of one fragment is exactly $\beta$$\cdot$N\_A, so its scheme spread is $\beta$ times the spread of the Becke population, 3.0415 kcal/mol per electron, i.e. 0.329 kcal/mol for the water dimer on the official S22 geometry and 0.300 on the constructed one (\texttt{nlcsplit/\allowbreak{}step\_\allowbreak{}k\_\allowbreak{}population.py}, whose two runs both assert that the fragment populations sum to the electron number).

\textbf{Attribution.} Four real-space partitions reachable through the PySCF grid machinery - Becke, Becke with the atomic-radius adjustment, Becke-Treutler and Stratmann - plus a nearest-atom hard assignment on a common grid. A point may be shared by more than one fragment; the factor matrix is exposed as \texttt{partition\_\allowbreak{}factors(mol, level, scheme) -> (coords, base\_\allowbreak{}w, S[N, natm])}, normalised so that \texttt{$\Sigma$\_a S[i,a] = 1} pointwise. The function refuses to return unless two internal assertions pass: normalisation, and \texttt{base\_w $\cdot$ S[i, owner(i)]} reproducing the weights that PySCF's own \texttt{get\_partition} produces - an external anchor against the same C kernel PySCF uses, not against a re-implementation in this package. The two checks it enforces are pointwise normalisation \texttt{|$\Sigma$\_a S[i,a] - 1| < 1e-12} and agreement with PySCF's own weights \texttt{max|base\_\allowbreak{}w$\cdot$S[i,owner(i)] - w\_\allowbreak{}ref| < 1e-11}, with a normalised mean deviation below 1e-14; the tolerances are the assertion thresholds in the published test, not residuals we are claiming to have measured \prov{S17}.

\textbf{Far-field cut.} Points below \texttt{rho\_floor = 1e-12} are dropped from the double sum. Sweeping 1e-8 \ldots{} 1e-16 moves the inter-fragment term by at most 2.1e-11 Ha at grid level 1 and 4.2e-11 Ha at level 2 - 1.3e-8 and 2.6e-8 kcal/mol, which the run's own log prints as 0.000000 - while discarding 10.7-22.6\% of points carrying $\leq$2.3e-6 electrons \prov{S7}, so the threshold is an implementation detail rather than a parameter. An earlier version of this sentence said "$\leq$5e-11 kcal/mol": it had taken the archived figure, which is in hartree, and put a kilocalorie unit on it, claiming the term was about 260 times less sensitive to the floor than the archive shows. Both figures are negligible against the 1e-4 kcal/mol scale of the quantities here, so the conclusion survives the error, but the number did not - corrected against \texttt{evidence/\allowbreak{}recompute/\allowbreak{}20260928T150222Z\_\allowbreak{}step\_\allowbreak{}d\_\allowbreak{}floor.out}, whose header already records the units.

\section{Software description}

\texttt{nlcsplit} is pure Python on top of PySCF's grid and density machinery, with NumPy for the pair sum; no compiled extension is built or required. Three modules carry the physics - \texttt{nlc.py} (VV10 fields and the block pair sum), \texttt{partition.py} (the six attribution schemes) and \texttt{geomlib.py} (geometries, fragment boundaries and the S22 loader). The \texttt{step\_*.py} scripts reproduce one table of the companion paper each and write a JSON sidecar; \texttt{figures.py} draws only from those sidecars, never from log text.

The public entry point is one function:

    res = nlc.fragment\_decomposition(coords, weights, rho, sigma, fragments,                                      b=b, C=C, owner=owner)      \# hard assignment     res = nlc.fragment\_decomposition(coords, weights, rho, sigma, fragments,                                      b=b, C=C, soft=S)           \# soft assignment

\texttt{fragments} is a list of atom-index groups. \texttt{owner} gives the nearest-atom index of each grid point; \texttt{soft} is an \texttt{(natm, ngrid)} weight matrix whose columns sum to one, as returned by \texttt{partition.scheme\_\allowbreak{}factors(mol, coords, scheme=...)}. The return value is a dict with \texttt{E\_loc}, \texttt{E\_nl}, \texttt{E\_total}, \texttt{intra} keyed by fragment, and \texttt{inter} keyed by fragment pair. \texttt{inter[(A, B)]} is the full cross contribution, not a half, so \texttt{E\_\allowbreak{}nl = sum(intra) + sum(inter)}; \texttt{closure\_check=True} compares the block sum against a monolithic evaluation over the same points and raises if they differ.

The six schemes are selected by name: the four Becke-family constructions used in the companion paper (\texttt{becke}, \texttt{becke-becke-rad}, \texttt{becke-treutler}, \texttt{stratmann}) and two cross-family controls (\texttt{voronoi}, \texttt{hirshfeld}). \texttt{scheme\_factors} returns the same shape for all six, so a comparison changes one keyword and nothing else, and it refuses to return a matrix whose columns do not sum to one within 1e-8.

A minimal run, condensed from \texttt{nlcsplit/\allowbreak{}examples/\allowbreak{}quickstart.py}:

    mol = gto.M(atom=geomlib.h2o\_dimer(), basis="6-31g*", unit="Angstrom")     mf, X0, W0, rho0, sig0 = partition.scf\_grid(mol, xc="wb97x\_v", level=1)     coords, weights, owner = partition.atom\_partition(mol, level=1, scheme="becke")     res = nlc.fragment\_decomposition(coords, weights, rho0, sig0,                                      geomlib.FRAGS["(H2O)2"], owner=owner)

For the water dimer at omegaB97X-V/6-31G*, grid level 1, this prints the numbers below, archived verbatim with its interpreter and module fingerprints as \texttt{evidence/\allowbreak{}release\_\allowbreak{}qa/\allowbreak{}20261003T1901Z\_\allowbreak{}quickstart\_\allowbreak{}readme\_\allowbreak{}example.out}:

    inter[(0,1)] = -0.5237     intra        = -3.5457 -3.7437     E\_nl         = -7.8130   E\_loc = 60.8303   E\_total = 53.0172

\texttt{E\_nl} equals \texttt{intra + intra + inter} as required, and \texttt{E\_total} may be compared with the native kernel on its own reference grid, which gives 53.017 and 53.016 for this system and level (difference +0.002 kcal/mol, the grid-representation difference of the companion text's section 3.1 rather than an algorithmic one) \prov{S9}.

Dependencies are NumPy and PySCF. \texttt{partition.py} reaches \texttt{libdft.VXCgen\_grid} through ctypes deliberately: it is the entry point \texttt{gen\_grid.get\_partition} itself uses, which makes the audit against PySCF's own partition external rather than circular. There is no other compiled dependency and no build step.

\section{Impact}

The package makes three things checkable that previously were not. It gives a number for how much of the nonlocal correlation is inter-fragment, which on the six S22 complexes that carry a counterpoise reference at grid level 2 accounts for 0.92-1.08 of the frozen-density superposition difference - association-time nonlocal correlation is almost entirely carried by point pairs straddling the fragments [S2, S10]. It supplies the per-system uncertainty that the attribution choice contributes (0.004-0.247 kcal/mol absolute, 0.8-18.2\% relative, with no ordering by interaction type or by system size), which is the error bar any future use of such numbers must carry \prov{S12}. That span is not basis-set invariant: on the water dimer it is 0.0249 kcal/mol in STO-3G against 0.0382-0.0400 in the four larger sets, so the number moves by \textasciitilde{}38\% between STO-3G and 6-31G* while the four sets beyond STO-3G agree with each other within 5\% \prov{S12}. Nor is it grid-converged to better than it is printed: level-2 values are stored to three decimals, so a level-1 to level-2 change below 5e-4 kcal/mol is not resolvable in the archived table, and the one four-decimal ladder we ran gives 0.0002 kcal/mol between levels 1 and 2 - small, but not zero. And it produces evidence about the quantity itself: where the share is defined at all it ranges from 16\% of $\Delta$E\_xc (water) to 97\% and 52\% for the two geometries of the benzene dimer, and on the methane dimer the two parts carry opposite signs so the package's own rule (|$\Delta$E\_NLC/$\Delta$E\_xc| > 2 $\Rightarrow$ report the two absolute terms, not a ratio) refuses to state a share at all. The VV10 term therefore cannot be reported as "the dispersion energy" without qualification, and on one row of our own table not even as a percentage \prov{S2}.

The per-system numbers behind these statements are tabulated in the companion paper; this paper does not repeat them.

\section{Availability and quality assurance}

\texttt{nlcsplit} is released under Apache-2.0. It contains no GAMESS source, and the claim is checked by scanning the published tree for source artefacts (\texttt{*.src}, \texttt{*.f}, \texttt{*.for}, \texttt{GAMESS*}) rather than for the word, and by installing the wheel into a throwaway target and walking what lands in \texttt{site-packages}. The archived scan reports 45 installed files, three named exemptions and no hit outside the exemption ledger at rc=0 \prov{S19}.

Correctness is guarded at three levels. The test suite pins the kernel against PySCF's native one, regression-tests the archived numbers, and checks the geometry and partition-factor invariants. The release script \texttt{tools/run\_ship\_qa.sh} enforces fifteen gates over the exported tree: build-residue and credential hygiene, no GAMESS source, byte-for-byte reproduction of every \texttt{evidence/*/MANIFEST.sha256}, no digit changed by the publishing scrub, import and wheel-member checks, and the provenance checks below. Each gate carries a two-directional self-test that must fire on a planted fault and stay silent on valid input, and the build refuses to ship when a self-test does not fire. The gates and the real defect that motivated each are listed in Table 1.

Two checks ship with the package itself. The ORCA 6.0.1 comparison recomputes the nonlocal term for the eight complexes in the companion paper, together with an omegaB97X control that must not print the label \prov{S18}. The partition-free referee re-evaluates the cross term on a two-Gaussian model density with three independent methods and reports the error bar \prov{S21}. Both are archived with their own manifests, and every \texttt{[S\#]} tag used in this paper resolves to a tracked file in the provenance appendix (Table 2); the converter re-checks each tag at build time and refuses to emit the PDF when one does not resolve.

The wheel census is part of the same script: members are counted by category (test modules, helpers under \texttt{nlcsplit/tests/}, everything else) and the category counts are printed beside their sum, so a silently dropped category cannot pass as a member count \prov{S19}.

\begin{table}[!htbp]
\centering
\footnotesize
\begin{tabularx}{\textwidth}{lXX}
\toprule
\textbf{gate} & \textbf{what it checks} & \textbf{standing} \\ \midrule
R & walks the \textbf{export tree} for build and pytest residue left inside it - \texttt{\_\_pycache\_\_} and \texttt{.pyc}, \texttt{.egg-info}, a stray \texttt{setup.cfg} - which is how a compiled docstring carries a repository path out even when every source file is clean & real defect: 13 \texttt{.pyc} files pytest had left inside the release directory, never opened by the old check, which iterated the list of copied names instead of walking the tree \\
A & no GAMESS source artefacts (matched on source patterns, not on the word "gamess", which the disclaimer itself contains) & real defect: a permanent false positive on our own disclaimer, which is why it matches patterns and not the word \\
B & walks the \textbf{published tree} and fails by default on credentials and internal-network strings - private-key blocks, password assignments, \texttt{ssh \ldots{} @} targets, bare-IP users - and, through the ten named rules that stay out of the release, on working-machine paths, bare cluster addresses and account names & real defect: the 13 published files that carried \texttt{/root/gamess\_work/...} and \texttt{/mnt/e/Project/GAMESS/...} while the gate reported every one of its arrays empty, because its patterns covered artefact extensions but not path forms \\
C & every \texttt{evidence/*/MANIFEST.sha256} reproduces byte-for-byte & real defect: a Windows path making \texttt{sha256sum} report six mismatches that were not there \\
C2 & the scrubbing changed no digit in any number, \textbf{and} every file in a published evidence directory is registered in that directory's own source manifest & the second half fired on real input the day it was written: a comparison JSON produced by a step script landed in \texttt{evidence/agentb\_sweep/} with no manifest entry, and the build refused to ship it (\texttt{shipped: false}) \\
D & every path cited by the companion text's provenance appendix is actually shipped & real defect: three citations that resolved to nothing in the published tree \\
E & every import in a shipped module is declared in \texttt{pyproject.toml} & real defect: \texttt{figures.py} importing matplotlib, which no dependency listed \\
F & the wheel and sdist member lists are what we think they are & real defect: two script-style tests that survived a build already reported fixed \\
G & \texttt{LICENSE} in the \textbf{source} file carries the real copyright line, not the Apache placeholder & real defect: the legal statement existed only in the generated copy; fires today on a restored placeholder \\
T & nothing in the export tree is missing from the export manifest & added as a belt-and-braces check; no firing on real input is recorded in this repository, only on the faults we planted \\
U & every module a shipped file imports from \texttt{nlcsplit} (absolute, dotted or package-relative) is itself in the shipped tree; the gate also self-tests against a planted missing sibling, and a build where that self-test does not fire does not ship & real defect: it was written for \texttt{fastpair.py}, which gate E could not see because an intra-package import is not a declared dependency \\
V & every Python module inside a packaged directory is either in the export whitelist or named, with a reason, as an unreleased internal module - this is the set \texttt{pip install .} actually puts in \texttt{site-packages}, which the whitelist does not constrain; self-tests fire on a planted stray and stay silent on the named ones & real defect: found by the shipped licence scanner, which walked 45 installed files from a source-directory install and found the peer implementation among them \\
W & no manuscript or README quotes a \emph{share} for a system whose denominator cancels, i.e. where \texttt{abs(dE\_NLC/dE\_xc) > 2} - the same rule \texttt{step\_f\_dispersion.py} applies to its own output, extended to the prose & real defect: three of them on the day it was written (the companion text's conclusion, this paper's section 4, and a line the author had just added to this README), all quoting a share equivalent to a ratio above 2 that section 3.2 of the same paper argues against - and the gate caught its own author reintroducing the number while documenting the gate \\
X & every \texttt{step\_*.py} in the export tree appears in \texttt{run\_all.sh}'s command table, so each table in the text has an entry point a reader can actually invoke & real defect: \texttt{step\_m\_ccsdt.py} reached the release on 2026-09-30 without ever being registered, which is the same omission the driver had been rewritten to fix - that rewrite left no guard, and this gate is the guard it lacked. The self-test truncates the command table to the four entries of the original bug and must then report ten missing scripts, and must go silent on a table completed by construction \\
Y & every image and link target in \texttt{paper/joss/paper.md} resolves to a file the export actually carries --- the gate walks the copied-file set, not the filesystem, so a figure that exists on the working machine but was never listed is a broken citation for the reader & planted only: it has not fired on real input, its two-directional self-test must fire on exactly the planted target and stay silent on a valid one, and it goes quiet when \texttt{paper/joss/} is absent from the tree \\
\bottomrule
\end{tabularx}
\caption{The release gates: what each checks, and the real defect that put it there. A letter identifies a gate in \texttt{tools/run\_ship\_qa.sh}; rows marked "planted only" have fired in self-test but not yet on real input.}
\label{tab:1}
\end{table}

\section{Limitations}

Quadratic cost in retained grid points, with no sub-quadratic approximation. One attempt at a certified screening was built and measured, and rejected - not on the weak version of the argument we first wrote here. The earlier sentence claimed the bound prunes nothing "at the tolerance the companion text needs (1e-8 Ha absolute)"; that mis-states both sides. \texttt{1e-8 Ha} is a per-node-pair acceptance constant, and the companion text needs 1e-6 kcal/mol on a fragment-pair term, i.e. 1e-6/627.509474 Ha in the same unit - so the constant and the requirement were never the same kind of quantity. Measured properly, three things hold: at atom-block granularity the \emph{smallest} block-pair bound in a real water dimer is 197.96 Ha against a quantity of -8.3e-4 Ha, so nothing is prunable at that granularity whatever the tolerance (0 of 21 pairs, and 0 of 300 on the 24-atom sandwich); refining below the atom block does fire, and at the accuracy the decomposition needs it still evaluates 84.4-95.4\% of the far-field point pairs; and with no other runnable user task at the sampling instant (a contention gate, not a claim that the machine was idle), the screened path is slower than the dense sum in every one of 33 timed screened/exact pairs across two independent windows on the same box - 1.08-10.28x in the archived window and 1.37-13.29x in the second (1.95-10.28x and 2.63-13.29x at the tightest tolerance tested), with no faster sample in either. The radius the bound subtracts is set by the far tail of each atom's density, so the certified separation comes out zero for every pair, including the ones facing each other across the inter-fragment gap. That is a statement about this particular bound, not about screening in general, and the measurements are in \texttt{paper/p2.md} with their artefacts under \texttt{evidence/bound\_forms/}, \texttt{evidence/fastpair\_ab/} and \texttt{evidence/scaling/} rather than quietly dropped \prov{S20}. Interpretation limited to the two-body and fragment-pair level. Results depend on the choice of real-space partition to the extent quantified in section 4. VV10 only - the $\tau$-dependent rVV10 form is not implemented. And it is not an energy decomposition analysis: it does not separate electrostatics, exchange-repulsion, polarisation or charge transfer.

\section{Appendix: sources}

Paths are relative to the repository root, because that is what the export gate resolves them against. Every path below is tracked by git, and the build refuses to produce \texttt{software.tex} if one is not. Three rows used to name \texttt{logs/<something>.out}, which is not a path in this repository: \texttt{nlcsplit/logs/*.out} is gitignored scratch, and those files are tracked only as the hash-manifested copies under \texttt{evidence/**} (the build does place a scrubbed copy of two of them inside the package, but a citation should point at one tracked original, not at a generated duplicate). The two \texttt{.json} logs in the S2 row are tracked in place, and are cited as they are kept.

Literature is not in this table. Where a sentence needs a published work it uses a \texttt{[[key]]} citation resolved against \texttt{paper/references.bib}, whose entries were reverse-checked through Crossref by DOI; the queries for the software and benchmark entries are archived in \texttt{evidence/\allowbreak{}bib\_\allowbreak{}provenance/\allowbreak{}G\_\allowbreak{}s22\_\allowbreak{}provenance\_\allowbreak{}20260930.md}.

\subsection{Artefact provenance}

\begin{table}[!htbp]
\centering
\footnotesize
\begin{tabularx}{\textwidth}{lXX}
\toprule
\textbf{Tag} & \textbf{Artefact} & \textbf{What it fixes} \\ \midrule
S1 & \texttt{RESEARCH-20260929-lit.md} \S{}1, \S{}6, \S{}9 & the literature screen behind the statement-of-need comparison, together with that file's own measurement that open-index \emph{subject} search recall is poor --- which is why \S{}1 calls the gap a screen and not an enumeration \\
S2 & \texttt{evidence/\allowbreak{}recompute/\allowbreak{}20260929T010500Z\_\allowbreak{}step1\_\allowbreak{}geomfix.out}, \texttt{evidence/\allowbreak{}recompute/\allowbreak{}MANIFEST.sha256}, \texttt{nlcsplit/logs/step\_h.json}, \texttt{nlcsplit/\allowbreak{}logs/\allowbreak{}step\_\allowbreak{}h\_\allowbreak{}l2cp.json}, \texttt{MAIN-skeleton.md} (its uncertainty-decomposition, P\_frz, and S22-subset tables) & the per-scheme block-sum residual - 1.73e-18 to 6.94e-18 Ha, printed for five partitions and two repeated schemes of the level-1 water dimer at lines 6-14 of the first artefact - the only measured additivity residual for a real system in this repository. The partition-free referee reports its own block-additivity residual on a two-Gaussian model density (worst 1.735e-18 Ha, row S21), which is a different object and not interchangeable with this one. Earlier versions of this row cited the two json logs for the residual: they carry the per-system tables and no additivity field at all \\
S3 & \texttt{nlcsplit/\allowbreak{}step\_\allowbreak{}f\_\allowbreak{}dispersion.py} & that \texttt{$\beta$$\int$$\rho$ = $\beta$N} cancels identically from a superposition difference, so the inter-fragment term needs no local-density convention --- and the denominator rule that refuses to quote a share when the sum cancels \\
S5 & \texttt{nlcsplit/step\_e\_kernel.py}, \texttt{evidence/\allowbreak{}recompute/\allowbreak{}20261001T083716Z\_\allowbreak{}step\_\allowbreak{}e\_\allowbreak{}pairing\_\allowbreak{}control.out}, \texttt{evidence/\allowbreak{}recompute/\allowbreak{}MANIFEST.sha256} & the definition-sensitivity control as both generator and record: the script compares same-centre with cross pairing on one grid and one density and prints the per-variant totals, and the archived stdout now carries the +0.1379/+0.4182/+0.1346 (level 1) and +0.9057 (benzene, level 0) kcal/mol of \S{}2. Until 2026-10-01 this row pointed at a re-measurement only, because the stdout that recorded them lived in nlcsplit/logs, which is gitignored - and the unarchived benzene figure in this paper (0.914) did not match the archived one in the companion text (0.906); the run resolved it in favour of 0.906 \\
S7 & \texttt{nlcsplit/step\_d\_floor.py}, \texttt{evidence/\allowbreak{}recompute/\allowbreak{}20260928T150222Z\_\allowbreak{}step\_\allowbreak{}d\_\allowbreak{}floor.out} & the \texttt{rho\_floor = 1e-12} far-field cut, and the reported cost of it (10.7--22.6\% of points dropped, carrying $\leq$ 2.3e-6 electrons) \\
S9 & \texttt{nlcsplit/\allowbreak{}step\_\allowbreak{}i\_\allowbreak{}gridrepr.py}, \texttt{evidence/\allowbreak{}recompute/\allowbreak{}20260929T\_\allowbreak{}step\_\allowbreak{}i\_\allowbreak{}gridrepr.out} & agreement of the explicit pair sum with PySCF's native VV10 kernel on a fixed grid, and the grid-representation residual quoted in \S{}1 and \S{}2 \\
S10 & \texttt{evidence/\allowbreak{}l2cp/\allowbreak{}20260928T213839Z\_\allowbreak{}centos-step\_\allowbreak{}h\_\allowbreak{}l2cp.out}, \texttt{evidence/\allowbreak{}l2cp/\allowbreak{}20260929T051912Z\_\allowbreak{}centos-l2\_\allowbreak{}benzmeth\_\allowbreak{}cp.out} & the level-2 counterpoise numbers, computed on the cluster and archived with their stdout \\
S12 & \texttt{evidence/\allowbreak{}recompute/\allowbreak{}20260929T003144Z\_\allowbreak{}stepg-recompute.out} & the basis-set and grid sensitivity of the attribution, including the $\pm$5\% agreement among the four basis sets beyond STO-3G \\
S17 & \texttt{nlcsplit/\allowbreak{}tests/\allowbreak{}test\_\allowbreak{}partition\_\allowbreak{}factors.py} & the two assertions \texttt{partition\_factors()} enforces before returning (pointwise normalisation at 1e-12, agreement with PySCF's own \texttt{get\_partition} weights at 1e-11) --- assertion thresholds, not measured residuals \\
S18 & \texttt{evidence/orca\_run/}, \texttt{evidence/orca\_control/}, \texttt{evidence/orca\_widen/}, \texttt{evidence/\allowbreak{}orca\_\allowbreak{}widen/\allowbreak{}cross\_\allowbreak{}program\_\allowbreak{}table.txt} & the ORCA 6.0.1 outputs and their manifests (\texttt{evidence/\allowbreak{}orca\_\allowbreak{}run/\allowbreak{}MANIFEST.sha256}, \texttt{evidence/\allowbreak{}orca\_\allowbreak{}control/\allowbreak{}MANIFEST.sha256}, \texttt{evidence/\allowbreak{}orca\_\allowbreak{}widen/\allowbreak{}MANIFEST.sha256}): six runs for the water and benzene-A anchors, the $\omega$B97X control group, and eighteen more covering the remaining six systems, with the archived stdout for every one of the 24; plus the two generators that re-derive the comparison from those outputs and refuse if the system list or any row's definition disagrees \\
S19 & \texttt{evidence/release\_qa/} & clean-environment installs and out-of-tree runs, each recorded with the SHA-256 of the wheel that was installed \\
S20 & \texttt{evidence/bound\_forms/}, \texttt{evidence/fastpair\_ab/}, \texttt{evidence/scaling/}, \texttt{evidence/coretail/} & the certified-screening measurements summarised in \S{}6, reported in full in \texttt{paper/p2.md} \\
S21 & \texttt{evidence/\allowbreak{}attribution\_\allowbreak{}referee/\allowbreak{}run.py}, \texttt{evidence/\allowbreak{}attribution\_\allowbreak{}referee/\allowbreak{}referee\_\allowbreak{}lib.py}, \texttt{evidence/\allowbreak{}attribution\_\allowbreak{}referee/\allowbreak{}summarize.py}, \texttt{evidence/\allowbreak{}attribution\_\allowbreak{}referee/\allowbreak{}README.md}, \texttt{evidence/\allowbreak{}attribution\_\allowbreak{}referee/\allowbreak{}outputs/\allowbreak{}summary.md}, \texttt{evidence/\allowbreak{}attribution\_\allowbreak{}referee/\allowbreak{}outputs/\allowbreak{}results.json}, \texttt{evidence/\allowbreak{}attribution\_\allowbreak{}referee/\allowbreak{}outputs/\allowbreak{}results\_\allowbreak{}run1\_\allowbreak{}20261001.json}, \texttt{evidence/\allowbreak{}attribution\_\allowbreak{}referee/\allowbreak{}outputs/\allowbreak{}table\_\allowbreak{}run1\_\allowbreak{}20261001.md}, \texttt{evidence/\allowbreak{}attribution\_\allowbreak{}referee/\allowbreak{}MANIFEST.sha256} & the partition-free referee for attribution: the two-Gaussian model density, the three independent evaluations of the exact cross term, the measured error bar it is compared against, and the two shipped runs whose tables are byte-identical apart from recorded wall-clock fields. Re-running \texttt{evidence/\allowbreak{}attribution\_\allowbreak{}referee/\allowbreak{}run.py} regenerates everything; \texttt{evidence/\allowbreak{}attribution\_\allowbreak{}referee/\allowbreak{}summarize.py} re-quotes it with no re-quadrature \\
\bottomrule
\end{tabularx}
\caption{Provenance appendix: what each [S\#] tag resolves to, and which claim in the text it fixes.}
\label{tab:2}
\end{table}
\bibliographystyle{unsrtnat}
\bibliography{references}
\end{document}
